\todo{Reprendre en concentrant sur SU}

%---------------------------------------------------------------------
\subsection*{Déroulement de carrière}
%---------------------------------------------------------------------
% Présentation chronologique des principales étapes de la carrière faisant apparaître les éléments les plus % significatifs (diplômes, positions, principales responsabilités et activités)

% %---------------------------------------------------------------------
% \paragraph{Formation initiale et titres universitaires.}
% \begin{itemize}
% \item[1988~:] ENSAE (\'Ecole Nationale de la Statistique et de l'Administration \'Economique, aujourd'hui ENSAI).
% \item[1989~:]  DEA de Statistique de l'université PARIS VI.
% \item[1993~:] Doctorat en Statistique de l'université PARIS V~: "Analyse de sensibilité de tests non-paramétriques adaptés aux données censurées" sous la direction de C. Huber (Université Paris V).
% \item[2002~:]  Habilitation à Diriger des Recherches en Mathématiques] de l'université d'\'Evry~: "Répartitions de motifs dans des séquences d'ADN".
% \end{itemize}

% %---------------------------------------------------------------------
% \paragraph{Enseignement supérieur et recherche.}
Ma carrière jusqu'à présent se découpe en trois grandes parties. À l'issue de ma thèse, j'ai d'abord été maître de conférences puis professeur (PR2 en 2003) à \APT de 1994 à 2003. En 2004, je suis devenu directeur de recherches (DR2, et DR1 en 2016) à l'INRAE jusqu'en 2021, date à laquelle j'ai rejoint \SU en tant que professeur (PR1). 

De 1994 à 2021, j'ai été rattaché à l'\UMRlong, que j'ai dirigée de 2004 à 2013, avant de rejoindre le Laboratoire de Probabilités, Statistique et Modélisation (LPSM) de \SU en 2021. 
J'ai effectué un détachement à l'INRAE de Jouy-en-Josas de 1999 à 2001, pendant lequel j'ai participé à la création de l'unité INRAE Mathématiques, Informatique et Génome (MIG). 
J'ai également été mis à disposition du Centre d'Écologie et des Sciences de la Conservation (CESCO) du Muséum National d'Histoire Naturelle entre 2020 et 2021.

%---------------------------------------------------------------------
\subsection*{\'Evolution des activités}
%---------------------------------------------------------------------
% Présentation de l’évolution éventuelles des activités

%---------------------------------------------------------------------
\paragraph{Enseignement.}
% APT
J'ai commencé ma carrière comme enseignant dans une école d'ingénieur, \APT, dont les étudiants sont pour une très large part issus de classes préparatoires BCPST (biologie, physique, chimie, sciences de la vie et de la Terre). Je me suis donc adressé très tôt à un public d’étudiants en sciences, mais dans des filières non principalement mathématiques. J’ai également eu très tôt à participer puis à concevoir des enseignements dans un cadre interdisciplinaire, en collaboration avec des collègues biologistes, agronomes ou écologues. 

Les statuts de maître de conférences et de professeur du ministère de l’Agriculture (auquel est rattachée \APT) sont calqués sur ceux de l’université. Durant cette période, j'ai donc assuré un service plein de 192h équivalent TD par an. Mes enseignements s’adressaient à tous les niveaux de l’école (1ère et 2ème année, année de spécialisation et master 2, école doctorale). Ils visaient à donner aux étudiants une formation statistique généraliste (statistiques descriptive, statistique inférentielle \cite{DRV99}, modèles linéaires, méthodes multivariés), mais certains enseignements étaient plus particulièrement associés à certaines spécialisations (planification expérimentale en agronomie, données censurées en santé humaine, modèle mixtes en sciences animales, modèles linéaires généralisés en écologie \cite{DBE15}). Je donnais également différents cours optionnels en probabilités (modélisation aléatoire, modèles markoviens, processus de naissances et morts). De 2004 à 2013, j'ai donné le cours de modèles linéaires dans le master 2 de Statistique de l'université d'Orsay.

\bigskip
% INRAE
J’ai poursuivi une partie de mes enseignements lorsque j’ai été recruté à l’INRAE. Notamment, j'ai participé à la création du master ‘Mathématiques pour les sciences du vivant’ (MathSV) de l’université Paris-Saclay/École polytechnique créé en 2012. L’objectif de ce master est de former de jeunes chercheuses et chercheurs en mathématiques appliquées intéressés par les sciences du vivant, en leur fournissant des bases solides en modélisation déterministe et en modélisation aléatoire. De 2012 à 2021, j’ai proposé dans ce master un cours sur les modèles à variables latentes orientés vers la biologie, la génomique, puis vers l’évolution et l’écologie. J’ai repris en grande partie ce cours en rejoignant \SU.

Pendant cette même période, j’ai continué à participer à un cours de modélisation statistique que j’avais co-créé pour l’école doctorale ABIES (agriculture, alimentation, biologie, environnement et santé). L’objectif de ce cours, réparti sur deux semaines pleines, est de former des doctorants en sciences du vivant à une démarche complète de modélisation, implémentation et analyse d’un problème issu de leurs recherches. Cette formation résultait du constat qu’il est difficile de former les étudiants à une telle démarche tant que les questions de recherche auxquelles elle doit permettre de répondre ne sont pas les leurs. 

Sur l'ensemble de cette période, j'ai assuré l'équivalent d'environ un quart de service d'enseignement universitaire par an.

\bigskip
Je décris mon activité d’enseignement à \SU dans la section \ref{sec:ens}. Ces enseignements se répartissent essentiellement en un cours L1 en Sciences des Données, un cours de L3 de Mathématiques pour les biologistes, d’un cours de M2 sur des modèles statistiques pour l’écologie et une série de cours dans un M2 de biologie évolutive et génétique des populations. Ces trois derniers enseignements bénéficient de mon expérience passée en enseignement et en recherche à l’interface entre mathématiques et sciences du vivant.

%---------------------------------------------------------------------
\paragraph{Recherche.}
Depuis le début de ma carrière, mes recherches se situent à l'interface entre les mathématiques (principalement les statistiques et les probabilités) et les sciences de la vie. La quasi-totalité de mes travaux plus méthodologiques ou théoriques ont été motivés, dans une plus ou moins grande mesure, par la biologie au sens large. Après avoir rédigé une thèse sur les tests non paramétriques utilisés en statistique médicale, je me suis rapidement intéressé aux applications en génomique, en particulier aux modèles probabilistes permettant de caractériser la distribution des motifs d'intérêt le long du génome \cite{RRS06}, qui a fait l'objet de mon habilitation à diriger des recherches (HDR). 

Au tournant des années 2000, j'ai accompagné l'arrivée des données transcriptomiques. En raison de leur nature massive, ces données ont contraint les biologistes moléculaires à recourir systématiquement et, si possible, judicieusement à la méthodologie statistique. J'ai alors proposé une série de contributions dans le domaine des tests multiples. À cette époque, j'ai également cofondé la conférence “Statistical Methods for Post Genomic Data” (\url{smpgd.fr/}), qui rassemble encore chaque année une centaine de chercheurs, principalement européens, pour discuter des questions liées aux derniers développements en biologie moléculaire.
Ces interactions avec la biologie m'ont amené à m'intéresser à trois thèmes généraux qui alimentent encore aujourd'hui mes recherches : ($a$) les modèles à variables latentes, ($b$) la détection des ruptures dans un signal et ($c$) les modèles de réseau. 
Ces trois familles de modèles présentent la particularité de poser des problèmes d’inférence statistique particuliers. 
Typiquement l’évaluation ou la maximisation de leur vraisemblance n’est généralement pas faisable par des techniques classiques. 
Il s’agit donc pour chacun d’entre eux de proposer des méthodes alternatives efficaces d’un point de vue computationnel et présentant des garanties statistiques suffisantes 
J’ai encadré, et encadre encore, plusieurs thèse sur chacun de ces sujets.
Afin de renouveler mes thématiques de recherche, j’ai décidé il y a une dizaine d’années de me tourner vers les applications en écologie. Cette transition s’est faite naturellement par le truchement de l’écologie microbienne qui se fonde sur des données essentiellement moléculaires. Tout comme pour la biologie moléculaire, j’ai découvert en m’en approchant l’immensité de ce champ disciplinaire dans lequel une culture statistique et une tradition de modélisation aléatoire sont bien installés. Par chance, ou du fait de leur caractère générique, les problématiques décrites plus haut se retrouvent en écologie, par exemple, ($a$) dans la modélisation de distributions jointes d’espèces, ($b$) dans l’analyse de comportements animaux ou ($c$) dans celle des réseaux plantes-pollinisateurs.

Mon arrivée au LPSM m’a permis de nouer de nouvelles collaborations autour de thématiques qui m’occupaient déjà. Elle m’a aussi permis de m'intéresser à de nouveaux objets comme les processus Hawkes (dont les applications en sciences du vivant sont légions) ou la statistique robuste.

% Les modèles à variables latentes ($a$) constituent une classe extrêmement large de modèles statistiques comme les modèles de mélange, les modèles mixtes ou les modèles de Markov cachés. Ils ont la particularité d'inclure des variables non observées (ou latentes), ce qui complique considérablement l'évaluation de leur vraisemblance. Depuis plusieurs années, une partie importante de mon travail consiste à étudier ces modèles et à proposer des algorithmes efficaces pour leur inférence. Initialement motivé par la détection des altérations génomiques locales rendue possible par des avancées technologiques, Je me suis également intéressé à la détection des ruptures ($b$). Du fait leur nature discontinue, ces modèles posent également des problèmes d'inférence intéressants, d'un point de vue combinatoire, ou en termes de sélection de modèles. Enfin, les réseaux constituent un formalisme naturel pour décrire les interactions entre différentes entités (gènes, molécules, espèces). La encore, des modèles statistiques à la fois interprétables et adaptés à ces observations sont nécessaires pour permettre de comprendre l’organisation d’un réseau. Le développement de modèles de graphes aléatoire hétérogène m’a ainsi à m’intéresser aux techniques d’approximations variationnelles qui ont ouvert de larges perspectives, tant méthodologiques que théoriques, qui alimentent encore aujourd’hui mes travaux.

Mes publications se répartissent, à part inégales et variables au cours du temps, entre des revues de statistique ou de probabilités appliquées canoniques, des revues d’interfaces (bioinformatique, biostatistique, écologie statistique) et, moins souvent, des revues de sciences du vivant (biologie moléculaire, écologie, évolution).

%---------------------------------------------------------------------
\paragraph{Responsabilités collectives et animation de la recherche.}
% Je donne le détail de mes responsabilités collectives à la section \ref{sec:col}. 
% La direction de l’équipe \SG de l’\UMRcourt a été ma première responsabilité collective en matière de recherche, il s’agissait d’une responsabilité essentiellement scientifique dans une équipe d’environ 10 membres travaillant sur des sujets proches. La responsabilité la plus importante est certainement la direction de l’\UMRcourt pendant 9 ans (trois évaluations de l’HCERES). L’unité de taille moyenne comptait alors un peu plus de 20 permanents et de 30 membres au total répartis en deux ou trois équipes selon les périodes et travaillant sur des sujets beaucoup plus variés en informatique, statistique et probabilités appliqués avec des applications en risque sanitaire, environnement ou génomique.  

% Pendant mon long séjour à \APT, j’ai toujours été élu dans une instance telle que le conseil des enseignants, les conseils scientifiques ou le conseil d’administration. J’ai également siégé dans les conseils scientifiques de plusieurs départements de l’INRAE, ainsi qu’au CNU, en section 26. J'ai enfin été membre du comité de pilotage de la fondation mathématique Jacques Hadamard (FMJH) et de l'école doctorale de mathématiques Hadamard (EDMH) lors de la création. Je suis ensuite devenu membre du conseil scientifique de la FMJH. Je n’ai assuré aucune responsabilité collective notable depuis mon arrivée à \SU, mais je suis membre du comité de pilotage de la fondation pour la science mathématique de Paris centre (FSMP) et du comité exécutif de l'intiative \SU InLife.

% Pour ce qui est de l’enseignement, comme indiqué plus haut, j’ai participé à la conception, la création et à la mise en place du master ‘Mathématiques pour les Sciences du Vivant” de l’université Paris-Saclay / École polytechnique et à la création et mise en place du master EvoGEM auquel sont associées toutes les universités parisiennes, ainsi que l’université Paris-Saclay.

%---------------------------------------------------------------------
% \subsection*{Formations suivies}
%---------------------------------------------------------------------
% Présentation des formations suivies notamment concernant vos activités pédagogiques
